\begin{tikzpicture}[
  >=Latex,
  x=1.25cm,
  y=1.2cm,
  line cap=round,
  line join=round,
  font=\small,
  ondaP/.style={->, blue!75!black, line width=1.6pt},
  ondaS/.style={->, orange!85!black, line width=1.6pt},
  ondaSup/.style={->, violet!75!black, line width=1.6pt}
]

% Corte esquemático del terreno.
\fill[gray!8] (-5.8,-1.85) rectangle (5.8,0.8);
\draw[black!75, line width=1.3pt] (-5.8,0.8) -- (5.8,0.8);
\node[above left, text=gray!65!black] at (-5.6,0.8) {superficie};

% Foco, epicentro y estación receptora.
\coordinate (foco) at (0,-1.3);
\coordinate (epicentro) at (0,0.8);
\coordinate (estacion) at (4.6,0.8);
\draw[densely dashed, gray!65] (foco) -- (epicentro);
\fill[red!75!black] (foco) circle (3pt);
\node[below left, align=right, text=red!70!black] at (-0.08,-1.3)
  {foco sísmico};
\fill[red!75!black] (epicentro) circle (2pt);
\node[above, text=red!70!black] at (0,0.88) {epicentro};

% Ondas internas P y S, y onda superficial.
\draw[ondaP] (foco) -- (estacion);
\draw[ondaS] (foco) .. controls (1.50,-1.10) and (3.10,-0.15) .. (estacion);
\draw[ondaSup] (0.18,0.81)
  .. controls (0.62,1.24) and (1.05,0.38) .. (1.48,0.81)
  .. controls (1.90,1.24) and (2.33,0.38) .. (2.76,0.81)
  .. controls (3.18,1.24) and (3.60,0.38) .. (4.18,0.81);

% Instrumento situado en la superficie.
\draw[black!70, line width=1pt] (4.6,0.8) -- (4.6,1.12);
\draw[black!70, line width=1pt] (4.42,1.12) -- (4.78,1.12);
\node[above right, align=left] at (4.82,1.38) {estación\\sísmica};

% Leyenda alineada bajo el corte.
\draw[ondaP] (-5.1,-2.3) -- (-4.35,-2.3);
\node[anchor=west, font=\footnotesize] at (-4.2,-2.3) {P: compresional};
\draw[ondaS] (-1.15,-2.3) -- (-0.4,-2.3);
\node[anchor=west, font=\footnotesize] at (-0.25,-2.3) {S: transversal};
\draw[ondaSup] (2.55,-2.3) -- (3.3,-2.3);
\node[anchor=west, font=\footnotesize] at (3.45,-2.3) {superficial};

\end{tikzpicture}